% Proposed history-latent replacement for the current DrifOv method.
% This formulation is not yet implemented in modeling_drif_ov.py.
% Suggested bibliography keys:
%   yang2026implicitdrifting  (arXiv:2606.01098)
%   black2025realtimechunking (arXiv:2506.07339)

\section{DrifOv: History-Latent One-Step Replanning}
\label{sec:drifov_method}

\subsection{Problem Formulation}

DrifOv targets asynchronous action-chunk generation for a robot that continues
to execute commands while the next policy request is in flight.  The method
must satisfy three requirements:
\begin{enumerate}
    \item generate the successor action chunk without iterative denoising;
    \item reproduce real-time chunking (RTC) during training, including its
          execution delay and immutable action overlap; and
    \item retain the executed action trajectory in a latent state so that the
          successor is conditioned on how the robot reached its predicted
          execution state.
\end{enumerate}

These requirements imply a separation between three types of action context:
\begin{equation}
    \underbrace{\mathcal{H}_t}_{\text{executed history}},
    \qquad
    \underbrace{U_{t:d}}_{\text{executed during inference}},
    \qquad
    \underbrace{P_{t+d}}_{\text{committed after arrival}}.
    \label{eq:three_action_contexts}
\end{equation}
The first describes the physical trajectory before the request.  The second
contains commands that move the robot while inference is running.  The third
contains commands that remain in the controller when the result arrives and
therefore cannot be replaced.  Treating these three sequences as a single
prefix loses their different causal roles.

The proposed model has a Y-shaped topology.  A causal history branch maintains
a latent representation of the executed trajectory.  A queue branch represents
the commands scheduled during and after inference.  The branches meet in a
one-step masked action decoder that predicts only the unknown suffix.  No
prefix classifier is required: committed actions are physical constraints, not
optional conditioning information.

\subsection{Asynchronous Execution Timeline}

Let an expert trajectory be
\begin{equation}
    \tau^n
    =
    \left\{
        (o_t^n,s_t^n,a_t^{n,*})
    \right\}_{t=0}^{T_n-1},
\end{equation}
where \(o_t\) is the visual-language observation, \(s_t\) is the robot state,
and \(a_t^*\in\mathbb{R}^{D}\) is the executed expert action.  The policy emits
a chunk of \(H\) actions.

Assume that a request is issued at control step \(t\).  Its estimated inference
delay is \(d\) control steps, so the result is intended to begin execution at
\begin{equation}
    r=t+d.
\end{equation}
At request time, the controller records the complete queue associated with the
request.  The commands expected to execute while inference is running are
\begin{equation}
    U_{t:d}
    =
    \left(
        a_t,
        a_{t+1},
        \ldots,
        a_{t+d-1}
    \right).
    \label{eq:inflight_actions}
\end{equation}
If \(K\) queued commands were available at request time and \(0\leq d\leq K\),
then \(L=K-d\) commands remain committed at arrival:
\begin{equation}
    P_{t+d}
    =
    \left(
        a_{t+d},
        a_{t+d+1},
        \ldots,
        a_{t+K-1}
    \right),
    \qquad
    L=K-d.
    \label{eq:committed_prefix}
\end{equation}
The delay-aligned training target begins at the intended arrival time:
\begin{equation}
    A_{t,d}^{*}
    =
    \left(
        a_{t+d}^{*},
        a_{t+d+1}^{*},
        \ldots,
        a_{t+d+H-1}^{*}
    \right).
    \label{eq:delay_aligned_target}
\end{equation}
For an aligned training sample, the first \(L\) target actions form the
committed prefix.  The remaining \(H-L\) actions are the continuation that the
model must generate.

Equations~\eqref{eq:inflight_actions}--\eqref{eq:delay_aligned_target} give the
delay a causal meaning.  Sampling \(d\) without shifting the target and
partitioning the queue does not reproduce asynchronous execution.

\paragraph{Action reference frame.}
All three action sequences must use a common execution-time reference frame.
Absolute joint commands satisfy this requirement directly.  Relative actions
must be re-anchored to the corresponding predicted or measured robot state.
Raw relative actions produced from different reference states must not be
copied or compared as if they were aligned.

\subsection{Y-Shaped Model Topology}
\label{sec:y_topology}

\subsubsection{Causal Historical Latent}

A single request-time observation cannot determine the arrival-time state when
the robot continues moving during inference.  DrifOv therefore maintains a
fixed-size latent memory
\begin{equation}
    Z_t
    =
    \left\{
        z_t^1,\ldots,z_t^M
    \right\},
    \qquad
    z_t^m\in\mathbb{R}^{C},
    \label{eq:latent_memory}
\end{equation}
consisting of \(M\) ordered latent tokens.  Multiple tokens are used instead of
one pooled vector so that the memory can retain different aspects of the
trajectory, such as scene state, robot configuration, recent motion, contact,
and task progress.

Let
\begin{equation}
    e_t
    =
    \left[
        E_o(o_t),
        E_s(s_t),
        E_a(a_{t-1}),
        E_{\Delta t}
    \right]
\end{equation}
be the evidence available at control step \(t\).  A causal memory update gives
\begin{equation}
    Z_t
    =
    M_\theta(Z_{t-1},e_t).
    \label{eq:memory_update}
\end{equation}
Only actions that were actually dispatched to the controller are written into
the historical memory.  Planned actions that are replaced before execution
must not be recorded as history.  In offline training,
Equation~\eqref{eq:memory_update} is unrolled over a finite history window.  During
deployment, the memory is updated online and cached between policy requests.

\subsubsection{Action-Conditioned Delay Propagation}

The latent \(Z_t\) represents the robot when the request is issued, whereas the
new chunk begins at \(t+d\).  A lightweight delay propagator advances the
historical latent using the known commands that will execute during inference:
\begin{equation}
    \bar Z_{t+d}
    =
    T_\theta
    \left(
        Z_t,
        E_U(U_{t:d}),
        E_d(d)
    \right).
    \label{eq:latent_rollout}
\end{equation}
Here, \(E_U\) encodes the complete in-flight sequence with execution-time
positions.  The propagator processes this sequence in one module evaluation;
it does not denoise or generate actions iteratively.  Its purpose is to predict
the latent control context at the intended arrival time.

The distinction between \(Z_t\) and \(\bar Z_{t+d}\) is essential.  A scalar
delay embedding tells the action decoder how much time has elapsed, but it does
not describe the state transition caused by the actions executed during that
time.  Equation~\eqref{eq:latent_rollout} conditions the transition on both.

\subsubsection{Committed-Prefix Tokens}

The commands remaining after arrival are represented independently from the
historical latent.  For output index \(h\in\{0,\ldots,H-1\}\), define the
coordinate-wise committed mask
\begin{equation}
    m_{h,j}
    =
    \mathbb{I}[h<L]\nu_{h,j},
    \label{eq:prefix_mask}
\end{equation}
where \(\nu\) is the valid-coordinate mask.  The action-query token is
\begin{equation}
    q_h
    =
    E_a(m_h\odot p_h)
    +E_m(m_h)
    +E_{\mathrm{exec}}(d+h)
    +E_{\mathrm{pos}}(h)
    +E_\xi((1-m_h)\odot\xi_h),
    \label{eq:action_query}
\end{equation}
where \(\xi_h\) is the source input used by the inherited one-step generator.
For deterministic endpoint regression, \(E_\xi\) may instead be a learned
unknown-action token.  Prefix coordinates are supplied explicitly, while
unknown coordinates contain no target-action information.

\subsubsection{One-Step Masked Action Decoder}

The decoder processes all \(H\) action queries jointly and cross-attends to the
arrival-time latent:
\begin{equation}
    \widetilde A_\theta
    =
    D_\theta
    \left(
        q_{0:H-1},
        \bar Z_{t+d}
    \right).
    \label{eq:one_step_decoder}
\end{equation}
This is one action-transformer evaluation.  The decoder may use bidirectional
self-attention across action queries because the complete output chunk is
generated simultaneously.  The history encoder remains causal, and the
decoder never receives the ground-truth suffix.

Committed actions are enforced after decoding:
\begin{equation}
    \widehat A_{h,j}
    =
    m_{h,j}p_{h,j}
    +(1-m_{h,j})\widetilde A_{\theta,h,j}.
    \label{eq:hard_prefix_copy}
\end{equation}
Consequently, the model cannot rewrite commands that the controller will
execute.  The committed prefix influences the generated suffix through
self-attention, but exact preservation does not depend on learned behavior.

\subsection{Training-Time RTC}
\label{sec:training_rtc}

Training examples are constructed from the same timeline used at deployment.
For each valid trajectory position:
\begin{enumerate}
    \item sample an inference delay \(d\) from the measured deployment latency
          distribution or a documented discrete approximation;
    \item sample a request-time queue length \(K\), with \(d\leq K<d+H\);
    \item encode history only through time \(t\) to obtain \(Z_t\);
    \item read \(U_{t:d}^{*}=a_{t:t+d-1}^{*}\), representing actions executed
          while the request is in flight;
    \item set \(L=K-d\) and construct the committed prefix
          \(P_{t+d}^{*}=a_{t+d:t+K-1}^{*}\);
    \item shift the target to \(A_{t,d}^{*}=a_{t+d:t+d+H-1}^{*}\);
    \item propagate \(Z_t\) through \(U_{t:d}^{*}\), decode once, and optimize
          only the unknown suffix together with latent-transition consistency.
\end{enumerate}

Samples crossing an episode boundary are masked or rejected.  The minibatch is
stratified over \(d\) and \(L\), including \(L=0\), so the decoder learns both
unconditioned generation and overlap-conditioned continuation.

\paragraph{No synthetic compatibility labels.}
The training prefix is not classified as valid or invalid.  A prefix taken from
a different time or perturbed in action space generally changes the physical
arrival state, making the original expert suffix an invalid counterfactual
target.  Model-generated prefixes may be used only when the corresponding
arrival observation, robot state, and continuation target are collected from a
simulator or real rollout.  Simple prefix corruption is therefore not part of
the core method.

\subsection{Learning Objectives}

\subsubsection{One-Step Suffix Generation}

Let
\begin{equation}
    \omega_{h,j}
    =
    \nu_{h,j}(1-m_{h,j})
\end{equation}
select coordinates for which the model is responsible.  We apply the inherited
one-step generation objective only to these coordinates:
\begin{equation}
    \mathcal{L}_{\mathrm{gen}}
    =
    \frac{
        \sum_{h,j}
        \omega_{h,j}
        \ell_{\mathrm{1step}}
        \left(
            \widetilde A_{\theta,h,j},
            A_{t,d,h,j}^{*}
        \right)
    }{
        \sum_{h,j}\omega_{h,j}
    }.
    \label{eq:generation_loss}
\end{equation}
For DrifOv, \(\ell_{\mathrm{1step}}\) can retain the one-step Drifting objective
of~\cite{yang2026implicitdrifting}; direct regression or a consistency-distilled
generator can use the same topology.  RTC conditioning does not require a
separate action, velocity, acceleration, attraction, or rejection loss.

\subsubsection{Arrival-Latent Consistency}

During training, the actual trajectory provides a representation of the
arrival-time history.  Applying the same causal memory encoder through
\(t+d\) gives a stop-gradient target
\begin{equation}
    Z_{t+d}^{+}
    =
    \operatorname{sg}
    \left[
        M_{\bar\theta}
        \left(
            \mathcal{H}_{t+d}
        \right)
    \right],
    \label{eq:arrival_latent_target}
\end{equation}
where \(\bar\theta\) denotes an exponential-moving-average copy of the memory
encoder or a stop-gradient online encoder.  The delay propagator is trained by
\begin{equation}
    \mathcal{L}_{\mathrm{dyn}}
    =
    \frac{1}{M}
    \sum_{m=1}^{M}
    \left\|
        \operatorname{norm}(\bar z_{t+d}^{m})
        -
        \operatorname{norm}(z_{t+d}^{+,m})
    \right\|_2^2.
    \label{eq:latent_dynamics_loss}
\end{equation}
This objective gives the historical latent a measurable role: after consuming
the in-flight actions, it must agree with the representation obtained from the
trajectory that actually reached the arrival time.

\subsubsection{Total Objective}

The core objective contains only generation and latent propagation:
\begin{equation}
    \mathcal{L}_{\mathrm{DrifOv}}
    =
    \mathcal{L}_{\mathrm{gen}}
    +
    \lambda_{\mathrm{dyn}}
    \mathcal{L}_{\mathrm{dyn}}.
    \label{eq:drifov_total}
\end{equation}
An auxiliary decoder from \(\bar Z_{t+d}\) to the arrival robot state may be
added if latent collapse is observed, but it is not a defining component.
Boundary smoothness is measured during evaluation and introduced as a loss
only if the base model exhibits a verified discontinuity that is not addressed
by prefix conditioning.

\subsection{Asynchronous Inference}

The deployment system maintains the latent memory and action queue as explicit
controller state:
\begin{enumerate}
    \item After each control step, update \(Z_t\) using the observation, robot
          state, and action that was actually executed.
    \item When issuing a policy request, snapshot \(Z_t\), estimate the arrival
          delay \(\widehat d\), and record the queue associated with the request.
    \item Partition that queue into the actions \(U_{t:\widehat d}\) expected to
          execute during inference and the prefix \(P_{t+\widehat d}\) expected
          to remain committed after arrival.
    \item Compute \(\bar Z_{t+\widehat d}\) with one delay-propagator evaluation
          and generate the complete candidate chunk with one action-decoder
          evaluation.
    \item At return time, compare the realized delay and queue state with the
          request metadata.  If the assumed overlap has expired or changed,
          discard the stale result rather than reinterpret it.
    \item Otherwise, preserve the committed prefix exactly and append only the
          newly generated suffix.
\end{enumerate}

The delay propagator does not perform iterative action generation.  The phrase
``one-step inference'' refers to replacing iterative diffusion or flow
integration with one action-transformer evaluation, consistent with the
one-step Drifting baseline.

\subsection{Information-Flow Constraints}

The topology uses architectural information boundaries rather than auxiliary
penalties:
\begin{enumerate}
    \item the historical encoder sees only observations, states, and executed
          actions at or before the request time;
    \item the delay propagator sees only the request-time latent, in-flight
          actions, and delay metadata;
    \item the action decoder sees the propagated latent, committed prefix,
          validity mask, timing positions, and one-step source input;
    \item the ground-truth suffix is used only by the loss;
    \item hard copying, rather than a learned penalty, preserves committed
          prefix coordinates.
\end{enumerate}
These constraints prevent future leakage during training and preserve the
causal interpretation of every input.

\subsection{Relation to RTC and One-Step Generation}

RTC~\cite{black2025realtimechunking} begins the next request before the current
chunk is exhausted and constrains the overlap during iterative diffusion or
flow inference.  DrifOv adopts the same asynchronous execution timeline during
training but replaces iterative inference-time guidance with three learned
components: causal trajectory memory, action-conditioned latent propagation,
and one-step masked continuation.  Exact prefix copying remains a controller
invariant.

The one-step generator itself is inherited from Drifting.  The new mechanism
addresses a different problem: a one-step policy must know not only the
request-time observation, but also the motion executed before its output can
take effect.  The history-latent branch supplies this missing control context.

\subsection{Implementation Changes}
\label{sec:implementation_changes}

The current \texttt{DrifOvActionHead} can be converted to the proposed topology
through the following changes:
\begin{enumerate}
    \item \textbf{Add persistent latent memory.}
          Maintain a fixed set of history tokens and update them with visual
          features, robot state, and the action actually executed.
    \item \textbf{Expose the request-time queue partition.}
          The RTC engine must provide both the actions expected to execute
          during inference and the commands expected to remain at arrival.
    \item \textbf{Add the delay propagator.}
          Encode the in-flight action sequence and transform the request-time
          memory into an arrival-time latent.
    \item \textbf{Replace global overlap metadata with execution-time tokens.}
          Each action query receives its committed mask and execution index
          \(d+h\).
    \item \textbf{Retain one-step decoding and hard prefix copying.}
          The existing action transformer and element-wise
          \texttt{torch.where} operation can remain the output path.
    \item \textbf{Replace boundary and compatibility objectives.}
          Train the unknown suffix with the inherited one-step objective and
          train the delay propagator with arrival-latent consistency.
    \item \textbf{Cache only physical history.}
          Update latent memory from actions confirmed as executed, not from the
          full action chunk originally predicted by the policy.
\end{enumerate}

\subsection{Required Ablations and Metrics}

The minimum ablation isolates each topological component:
\begin{enumerate}
    \item unconditioned one-step Drifting;
    \item prefix-conditioned generation without delay alignment;
    \item delay-aligned RTC with prefix conditioning but no historical latent;
    \item historical latent without action-conditioned delay propagation;
    \item the complete history-latent DrifOv model;
    \item complete model without \(\mathcal{L}_{\mathrm{dyn}}\);
    \item recurrent latent tokens versus a single pooled history vector.
\end{enumerate}

In addition to task success, evaluation reports:
\begin{enumerate}
    \item suffix action error and success rate by delay \(d\) and overlap \(L\);
    \item arrival-state or arrival-feature prediction error;
    \item exact committed-prefix preservation error;
    \item action discontinuity at the prefix--suffix boundary;
    \item stale-result frequency and achieved control update rate;
    \item history-memory update cost, delay-propagator cost, and action-decoder
          latency separately.
\end{enumerate}

This experiment distinguishes gains from one-step generation, training-time
RTC, and historical trajectory modeling instead of combining them into a
single collection of auxiliary constraints.
